\chapter{Stato dell'implementazione}
\label{ch:stato-implementazione}

Il presente capitolo documenta lo stato attuale del prototipo operativo \textit{Forecasting}, sviluppato in parallelo alla proposta descritta nel Capitolo~\ref{ch:descrizione-progetto}.
L'obiettivo è rendere esplicito cosa è già funzionante, con quale architettura e quali limiti, distinguendo le componenti implementate da quelle ancora pianificate.

\section{Architettura e stack tecnologico}

Il codice sorgente è organizzato in un repository Git e implementato interamente in \textbf{Python}.
L'applicazione è containerizzata con Docker e composta da tre elementi principali:

\begin{itemize}
  \item \textbf{PostgreSQL} per la persistenza strutturata degli articoli, dei prezzi e dei metadati aziendali;
  \item \textbf{Flask} (\texttt{app.py}) come server web e API REST;
  \item \textbf{Gunicorn} per il deploy in produzione su Heroku.
\end{itemize}

La configurazione locale avviene tramite \texttt{docker-compose.yml}; in produzione l'applicazione è esposta pubblicamente su Heroku all'indirizzo:
\url{https://dollarpunk-26786f215068.herokuapp.com/dashboard}

Rispetto al disegno iniziale, non sono stati ancora introdotti moduli in Rust o JavaScript dedicati, né pipeline di training su infrastrutture GPU (RunPod, Databricks).
L'approccio attuale privilegia un \textit{MVP} orientato all'ingestion, alla persistenza e alla visualizzazione.

\section{Raccolta dati}

\subsection{Fonte attuale: Alpha Vantage}

Al momento la raccolta testuale si appoggia esclusivamente all'endpoint \texttt{NEWS\_SENTIMENT} di Alpha Vantage, orchestrato dal modulo \texttt{news\_collector.py}.
Per ogni articolo l'API restituisce metadati (titolo, URL, fonte, timestamp, summary) e, per ciascun ticker menzionato, un punteggio di \emph{relevance} e di \emph{sentiment} già calcolati dal provider.

La struttura salvata in database replica il formato JSON dell'API, coerente con l'output atteso nel Capitolo~\ref{ch:descrizione-progetto}:

\begin{itemize}
  \item \texttt{stocks\_mentioned} e \texttt{ticker\_sentiments} (campo JSONB \texttt{ticker\_sentiment});
  \item \texttt{relevance\_score} e \texttt{sentiment\_score} per ticker;
  \item etichetta qualitativa (\texttt{Bullish}, \texttt{Bearish}, \texttt{Neutral}, ecc.).
\end{itemize}

\subsection{Modalità di ingestion}

Sono disponibili più modalità di acquisizione, attivabili dalla dashboard amministrativa o da job schedulati:

\begin{itemize}
  \item \textbf{Collect manuale}: raccolta parametrizzata per ticker, topic e intervallo temporale;
  \item \textbf{Deep ingestion / deep research}: acquisizione estesa su ticker o topic con tracciamento dello stato in tempo reale;
  \item \textbf{Coverage ingestion} (\texttt{coverage\_ingestion\_job.py}): scansione sistematica degli ultimi $N$ giorni su tutti i topic supportati dall'API e sui ticker presenti nel database, con progresso salvato su file e ripresa automatica;
  \item \textbf{Prune old news} (\texttt{prune\_old\_news.py}): manutenzione del database per mantenere una finestra temporale rolling degli articoli.
\end{itemize}

Tutte le esecuzioni significative sono registrate nella tabella \texttt{job\_executions}, consultabile dalla pagina \textit{Job Executions}.

\subsection{Scostamento rispetto al disegno iniziale}

Non sono ancora integrate le fonti eterogenee descritte nel Capitolo~\ref{ch:descrizione-progetto} (Bloomberg, Reuters, social media, Reddit, scraping con \texttt{newspaper}, motori di ricerca).
La raccolta attuale è quindi \textbf{monofonte} e dipendente dai limiti di rate dell'API Alpha Vantage (5 chiamate/minuto nel tier gratuito, fino a 75/minuto nel tier premium).

\section{Persistenza e modello dati}

Il modello relazionale principale è definito in \texttt{schema.sql}.
La tabella centrale \texttt{articles} memorizza ogni notizia con inserimento \textbf{idempotente} sulla chiave \texttt{url}.

Oltre agli articoli, il database include:

\begin{itemize}
  \item \texttt{stock\_prices}: prezzi di chiusura e volumi giornalieri per ticker, recuperati lazy da Alpha Vantage e cachati localmente;
  \item \texttt{companies}: anagrafica e descrizione del business per ticker (Alpha Vantage OVERVIEW e, opzionalmente, FMP);
  \item \texttt{company\_embeddings}: vettori embedding OpenAI (\texttt{text-embedding-3-small}) delle descrizioni aziendali;
  \item \texttt{company\_correlation\_matrix}: matrice di similarità coseno tra embedding aziendali, con cache in database;
  \item \texttt{job\_executions}: log strutturato dei job di background.
\end{itemize}

Sono inoltre presenti tabelle predisposte per estensioni future (\texttt{ticker\_decayed\_sentiment}, \texttt{ticker\_cross\_diffused\_sentiment}) la cui logica di popolamento non è ancora implementata.

\subsection{Viste analitiche SQL}

Per l'aggregazione del sentiment a livello di ticker sono state definite due viste:

\begin{itemize}
  \item \texttt{article\_ticker\_sentiment\_view}: espande il JSONB \texttt{ticker\_sentiment} in una riga per coppia (articolo, ticker);
  \item \texttt{ticker\_daily\_sentiment\_view}: aggrega per ticker e giorno, con decadimento temporale del \emph{relevance score} e media ponderata del sentiment.
\end{itemize}

\subsection{Archiviazione cold-storage}
\label{sec:cold-storage}

La finestra rolling di 30 giorni, imposta dai limiti del piano PostgreSQL in produzione (1~GB), comportava originariamente la perdita definitiva degli articoli rimossi dal job di prune.
Per eliminare questa perdita informativa è stata introdotta un'architettura a due livelli:

\begin{itemize}
  \item \textbf{Hot store}: il PostgreSQL di produzione, limitato agli ultimi 30 giorni, che serve dashboard e API con tempi di risposta contenuti;
  \item \textbf{Cold store}: un archivio analitico \textbf{DuckDB} ospitato su MotherDuck, alimentato quotidianamente e privo di limiti pratici di crescita.
\end{itemize}

Il job di archiviazione (\texttt{archive\_job.py}) copia articoli e prezzi giornalieri dal database operativo all'archivio con \textbf{inserimento idempotente} tramite anti-join sulla chiave naturale (\texttt{url} per gli articoli, coppia ticker--data per i prezzi): l'esecuzione ripetuta non produce duplicati.
Nello scheduler di produzione il job precede il prune, garantendo che nessun articolo venga cancellato prima di essere stato archiviato.

Al momento del bootstrap (luglio 2026) l'archivio conteneva \textbf{758\,817 articoli} con scoring per ticker --- con storico dal novembre 2018, recuperato in parte da un'istanza locale del database --- e 24\,758 osservazioni di prezzo.
Questo dataset costituisce la base empirica per due componenti della proposta originaria rimaste finora teoriche: il \emph{fine-tuning} di modelli open-source su coppie (testo, sentiment per ticker) già etichettate, e il backtesting su orizzonti pluriennali richiesto dal protocollo anti-bias del Capitolo~\ref{ch:direzioni-future}.

\section{Scoring del sentiment: delega al provider}

Una distinzione importante rispetto al Capitolo~\ref{ch:descrizione-progetto}: al momento \textbf{non esiste una pipeline NLP proprietaria}.
L'estrazione dei ticker, il \emph{relevance score} e il \emph{sentiment score} sono calcolati interamente da Alpha Vantage e persistiti così come ricevuti.

Di conseguenza, non sono ancora stati implementati:

\begin{itemize}
  \item modelli FinBERT, GPT-4 o LLM open-source per scoring interno;
  \item approcci RAG (FAISS, LangChain);
  \item fine-tuning su dataset annotati;
  \item NER, pattern matching e \emph{confidence score} proprietari per l'associazione testo--ticker.
\end{itemize}

Questa scelta consente di validare rapidamente il flusso ETL, lo storage e le visualizzazioni, rimandando lo sviluppo del modulo NLP a una fase successiva.

\section{Decay temporale e KPI di sentiment}

Il decadimento temporale descritto nel Capitolo~\ref{ch:descrizione-progetto} è stato implementato in forma operativa, con una parametrizzazione basata su \textbf{emivita} (\textit{half-life}) anziché sul parametro $\lambda$ esplicito del modello teorico.

Per ogni ticker, il KPI giornaliero calcolato in dashboard e nelle API pubbliche è:

\[
\text{KPI}_d = \frac{\sum_i s_i \cdot r_i \cdot w_i}{\sum_i r_i \cdot w_i},
\qquad
w_i = 0{,}5^{\,\Delta t_i / h}
\]

dove $s_i$ e $r_i$ sono rispettivamente \emph{sentiment score} e \emph{relevance score} dell'articolo $i$, $\Delta t_i$ i giorni trascorsi dalla pubblicazione e $h$ l'emivita configurabile (default: 7 giorni).
Il fattore $0{,}5^{\Delta t / h}$ equivale a un decadimento esponenziale con $\lambda = \ln(2)/h$.

La stessa logica è replicata nell'endpoint \texttt{/api/tickers/<ticker>/sentiment-kpi} e nella pagina portfolio (sentiment decayed su intervalli configurabili).
La vista SQL \texttt{ticker\_daily\_sentiment\_view} (lookback 30 giorni) adotta invece la variante \emph{non normalizzata}: una somma dei contributi giornalieri con decadimento, senza divisione per la somma dei pesi.
Le due grandezze non sono intercambiabili: come discusso nel Capitolo~\ref{ch:descrizione-progetto}, la media normalizzata misura la direzione del sentiment attivo ed è confinata nell'intervallo dei sentiment osservati, mentre la somma decayed ne misura l'intensità e tende a zero in assenza di notizie.
La coesistenza delle due varianti è quindi una scelta informativa, non una duplicazione.

\subsection{Diffusione semantica}

La proposta di diffusione del \emph{relevance} nello spazio semantico (equazione alle derivate parziali del Capitolo~\ref{ch:descrizione-progetto}) non è ancora operativa sulle notizie.
Esiste tuttavia un \textbf{primo passo parziale} sulle aziende: le descrizioni di business vengono vettorizzate con OpenAI e confrontate tramite similarità coseno, producendo una matrice di correlazione tra ticker consultabile dalla sezione \textit{Companies}.
Questa matrice misura affinità tra profili aziendali, non tra articoli di news, e non alimenta ancora la tabella \texttt{ticker\_cross\_diffused\_sentiment}.

\section{Indicatori di analisi tecnica}
\label{sec:analisi-tecnica}

Accanto al KPI di sentiment, il prototipo calcola un insieme di indicatori tecnici classici a partire dalle chiusure giornaliere già persistite in \texttt{stock\_prices}, senza chiamate API aggiuntive.
L'obiettivo è affiancare al segnale informativo (``cosa dicono le notizie'') un segnale di prezzo (``cosa sta facendo il titolo''), abilitando la lettura congiunta dei due.
Benché l'analisi tecnica sia storicamente controversa, la sua fondazione statistica è stata riabilitata in letteratura: \textcite{lo2000foundations} mostrano su dati di decenni che i pattern tecnici contengono informazione incrementale, e \textcite{neely2014forecasting} che gli indicatori tecnici prevedono il premio al rischio azionario con potere comparabile---e complementare---alle variabili macroeconomiche.

\begin{itemize}
  \item \textbf{Medie mobili semplici} a 20 e 50 giorni, con la distanza percentuale del prezzo dalla SMA~50 come indicatore sintetico di momentum;
  \item \textbf{Bande di Bollinger} \parencite{bollinger2001bollinger}: $\text{SMA}_{20} \pm 2\sigma_{20}$, dove $\sigma_{20}$ è la deviazione standard mobile a 20 giorni; la posizione del prezzo all'interno della banda, normalizzata in $[0,1]$, segnala movimenti anomali rispetto alla volatilità recente;
  \item \textbf{RSI a 14 giorni} \parencite{wilder1978new}, con smoothing di Wilder, e le soglie convenzionali di ipercomprato ($>70$) e ipervenduto ($<30$);
  \item \textbf{Volatilità realizzata a 20 giorni}, deviazione standard dei log-rendimenti giornalieri annualizzata con fattore $\sqrt{252}$.
\end{itemize}

Gli indicatori sono esposti da un endpoint dedicato (\texttt{/api/tickers/<ticker>/technicals}) e visualizzati in due punti dell'interfaccia: nella pagina di dettaglio del ticker, come overlay attivabili sul grafico prezzo--volume con uno snapshot dei valori correnti, e nella dashboard pubblica, come badge sintetici (direzione rispetto alla SMA~50, RSI) accanto al KPI di sentiment di ciascun titolo in classifica.

Sopra questi valori è costruito un \textbf{report tecnico automatico}: un generatore deterministico a regole che traduce gli indicatori in una lettura testuale (momentum, incroci di medie mobili, condizioni di ipercomprato/ipervenduto, \textit{squeeze} di volatilità) e la conclude con un confronto esplicito tra la direzione tecnica e il KPI di sentiment corrente, segnalando l'\emph{allineamento} o la \emph{divergenza} dei due segnali.
La rilevazione automatica delle divergenze (notizie bullish su un titolo tecnicamente debole, o viceversa) rende osservabile per ogni titolo, in tempo reale, il fenomeno al centro della domanda di ricerca del Capitolo~\ref{ch:direzioni-future}.

La co-presenza dei due segnali sulla stessa vista non è solo una scelta di interfaccia: prepara la domanda di ricerca sull'\emph{interazione} tra sentiment e segnali tecnici discussa nel Capitolo~\ref{ch:direzioni-future}, in linea con la letteratura che documenta il valore predittivo congiunto di regole tecniche e variabili informative \parencite{brock1992simple}.

\section{Interfaccia web e funzionalità utente}

L'applicazione Flask espone due livelli di accesso:

\paragraph{Area amministrativa} (autenticata)
Pannello di controllo per lanciare job di ingestion, consultare statistiche globali, esplorare articoli, eseguire query SQL e gestire l'anagrafica aziende con vectorizzazione batch.

\paragraph{Area pubblica} (read-only)
Dashboard giornaliera (\texttt{/dashboard}) con ranking dei ticker per KPI di sentiment decayed, filtri per direzione (bullish/bearish) e finestra temporale.
Per ogni ticker è disponibile una pagina di dettaglio con:

\begin{itemize}
  \item grafico prezzo e volume (Alpha Vantage, tabella \texttt{stock\_prices});
  \item serie temporale del KPI di sentiment;
  \item elenco delle ultime notizie che menzionano il ticker;
  \item metadati aziendali (settore, industry, descrizione).
\end{itemize}

La pagina \textit{Portfolio} consente di monitorare una watchlist personalizzata (persistita lato client) con quote, storico prezzi, sentiment decayed e ultime news per i ticker selezionati.

\section{Confronto sintetico: previsto vs.\ implementato}

La Tabella~\ref{tab:previsto-implementato} riassume lo scostamento tra la proposta del Capitolo~\ref{ch:descrizione-progetto} e lo stato corrente del codice.

\begin{table}[H]
\centering
\renewcommand{\arraystretch}{1.3}
\begin{tabular}{|p{5.5cm}|p{4cm}|p{4.5cm}|}
\hline
\textbf{Componente} & \textbf{Previsto} & \textbf{Stato attuale} \\
\hline
Fonti dati & Multi-fonte (news, social, scraping) & Solo Alpha Vantage \\
\hline
NLP / scoring & FinBERT, RAG, fine-tuning & Delegato al provider API \\
\hline
ETL e storage & Pipeline strutturata & Implementato (PostgreSQL) \\
\hline
Archiviazione storico & --- (non prevista) & Cold store DuckDB/MotherDuck, 758k articoli dal 2018 \\
\hline
Indicatori tecnici & Input per LSTM & SMA, Bollinger, RSI, volatilità in dashboard e API \\
\hline
Decay temporale & $R \cdot e^{-\lambda \Delta t}$ & Emivita $0{,}5^{\Delta t/h}$ in SQL e API \\
\hline
Diffusione semantica & PDE su spazio news & Solo embedding aziende (parziale) \\
\hline
Analisi sentiment--prezzo & Correlazioni Pearson/Spearman & Grafici separati, no statistica \\
\hline
Backtesting & Strategia long/short, Sharpe & Non implementato \\
\hline
Simulazione portafoglio & Capitale, costi, stop-loss & Watchlist + monitoraggio \\
\hline
LSTM predittivo & Prezzo + sentiment & Non implementato \\
\hline
Deploy & Docker, PaaS & Docker + Heroku operativi \\
\hline
\end{tabular}
\caption{Confronto tra componenti previste e stato di implementazione al momento della stesura.}
\label{tab:previsto-implementato}
\end{table}

\section{Prossimi passi}

Sulla base del prototipo attuale, le priorità di sviluppo coerenti con la proposta di ricerca sono:

\begin{enumerate}
  \item introdurre almeno una fonte dati aggiuntiva (scraping o social) per ridurre la dipendenza da un unico provider;
  \item implementare un modulo NLP interno (es.\ FinBERT zero-shot) per confronto con i punteggi Alpha Vantage, sfruttando il dataset etichettato ormai accumulato nel cold store (Sezione~\ref{sec:cold-storage});
  \item popolare e validare le tabelle di \emph{cross-diffused sentiment} sfruttando la matrice di correlazione tra aziende;
  \item aggiungere analisi esplorativa sentiment--prezzo con correlazioni e ritardi temporali, estesa all'interazione con gli indicatori tecnici (Sezione~\ref{sec:analisi-tecnica});
  \item costruire un backtester elementare prima di passare a modelli LSTM.
\end{enumerate}

Il prototipo dimostra la fattibilità della catena \emph{raccolta $\rightarrow$ persistenza $\rightarrow$ decay $\rightarrow$ visualizzazione}, che costituisce il substrato su cui innestare le componenti analitiche e predittive descritte nel capitolo precedente.
Le priorità qui elencate sono sviluppate in dettaglio, insieme alle corrispondenti domande di ricerca, nel Capitolo~\ref{ch:direzioni-future}.
